\section{Arquitectura General del Sistema}
El sistema se diseña bajo una arquitectura modular compuesta por dos componentes principales: el 
\textbf{componente de extracción} y el \textbf{componente de visualización}.
\medskip

El \textbf{componente de extracción} se encarga de cargar y ejecutar el modelo convolucional YOLOv2, 
capturar las activaciones internas y transformar dicha información a un formato de intercambio intermedio. 
Específicamente, este componente se estructura en tres módulos funcionales:

\begin{itemize}
    \item \textbf{Módulo de ejecución:} Gestiona la carga de la arquitectura y los pesos del modelo YOLOv2, 
    controlando la ejecución de la inferencia sobre las imágenes de entrada.
    \item \textbf{Módulo de extracción:} Configura e intercepta los mecanismos de captura (\textit{hooks}) 
    para extraer los mapas de activación, tensores intermedios y metadatos relevantes durante el procesamiento 
    de la red.
    \item \textbf{Módulo de formato e intercambio:} Procesa y empaqueta las matrices numéricas hacia estructuras 
    de datos estandarizadas (archivos JSON para metadatos y volúmenes OpenVDB para tensores 3D), facilitando la 
    persistencia y el transporte de información entre componentes.
\end{itemize}
\medskip

Por su parte, el \textbf{componente de visualización} se consolida como un complemento (\textit{add-on}) 
integrado en Blender. Su función es consumir los datos del formato intermedio, generar las representaciones 
volumétricas correspondientes y proveer una interfaz interactiva para su exploración. Este componente consta 
de los siguientes módulos:

\begin{itemize}
    \item \textbf{Módulo de lectura y preparación (\textit{Parsing}):} Carga e interpreta las estructuras de 
    datos procedentes del formato intermedio (JSON y OpenVDB), deserializando y preparando los tensores para 
    su manipulación dentro del motor gráfico.
    \item \textbf{Módulo de mapeo y representación visual:} Traduce proceduralmente los atributos volumétricos 
    y metadatos en entidades geométricas, materiales y nodos dentro de Blender, habilitando la inspección e 
    interacción tridimensional con el modelo.
\end{itemize}

Cada componente se desarrollará y ejecutará en un entorno independiente. Esta decisión de diseño 
responde a las siguientes consideraciones técnicas:

\begin{itemize}
    \item \textbf{Gestión de recursos y estabilidad:} Aunque técnicamente es factible instalar y ejecutar 
    modelos de inteligencia artificial dentro del intérprete de Python integrado en Blender, la ejecución 
    simultánea de tareas computacionalmente intensivas (como el cálculo matricial de inferencia y el renderizado 
    gráfico) puede saturar la VRAM/RAM y comprometer la estabilidad del sistema operativo.
    \item \textbf{Aislamiento de entornos y viabilidad técnica:} En entornos de producción e investigación, 
    la ejecución de modelos de IA se realiza habitualmente en hardware especializado y en entornos controlados. 
    Acoplar la inferencia del modelo y el motor de Blender dentro de un mismo ecosistema de ejecución introduce 
    fricciones de compatibilidad de dependencias y complejidad técnica innecesaria.
    \item \textbf{Modularidad y extensibilidad del software:} La separación de responsabilidades entre la 
    extracción y la visualización simplifica el desarrollo y depuración de ambos componentes. Esta arquitectura 
    desacoplada permite extender, actualizar o sustituir módulos de forma independiente 
    (por ejemplo, adaptar el extractor para otro modelo de visión sin modificar el complemento de Blender).
    \item \textbf{Gestión de artefactos y reproductibilidad:} El uso de un formato de intercambio intermedio 
    facilita la persistencia y administración de las activaciones extraídas. Esto permite a los usuarios 
    auditar inferencias previas sin necesidad de reejecutar el modelo con los parámetros originales, 
    funcionando como un registro histórico de versiones para evaluar la evolución del algoritmo. Asimismo, 
    habilita la portabilidad de datos entre colaboradores, permitiendo la exploración sin requerir 
    la ejecución local de la red de IA.
\end{itemize}

\begin{figure}[!htbp]
    \centering
    \includegraphics[width=0.8\textwidth]{diagrams/chapter3/general_architecture.pdf} 
    \caption[Arquitectura general del sistema]{Diagrama de la arquitectura general del sistema.}
    \label{fig:arquitectura_general_sistema}
\end{figure}